% Generated draft: requires mathematical and source review.
\documentclass[11pt]{article}
\usepackage[margin=25mm]{geometry}
\usepackage{fontspec}
\setmainfont{Latin Modern Roman}
\usepackage{amsmath,amssymb,mathtools,unicode-math}
\setmathfont{Latin Modern Math}
\usepackage{enumitem,needspace,xurl,hyperref}
\setlength{\parindent}{0pt}
\setlength{\parskip}{5pt}
\setlength{\emergencystretch}{3em}
\newcommand{\N}{\mathbb{N}}
\newcommand{\Z}{\mathbb{Z}}
\newcommand{\Q}{\mathbb{Q}}
\newcommand{\R}{\mathbb{R}}
\newcommand{\C}{\mathbb{C}}
\newcommand{\F}{\mathbb{F}}
\newcommand{\A}{\mathbb{A}}
\newcommand{\PP}{\mathbb{P}}
\newcommand{\E}{\mathbb{E}}
\newcommand{\eps}{\varepsilon}
\newcommand{\dd}{\,\mathrm{d}}
\begin{document}
\section*{UnsolvedMath ID: 800890}
\textbf{Source problem code:} catalogue entry
\subsection{Definitions and mathematical statement}
Let \((X,\omega)\) be a compact K\"ahler manifold of complex dimension \(n\), and let \(D\subset X\) be a divisor, with \(X\setminus D\) understood as the complement of its support. Fix the convention \(dd^c=\frac{i}{\pi}\partial\bar\partial\). A function \(u\colon X\to[-\infty,\infty)\) is \(\omega\)-plurisubharmonic if it is locally the sum of a plurisubharmonic function and a smooth function, is locally integrable, and satisfies \(\omega+dd^c u\geq 0\) as a current. Its non-pluripolar Monge--Amp\`ere measure is
\[
\left\langle(\omega+dd^c u)^n\right\rangle
=\lim_{j\to\infty}\mathbf{1}_{\{u>-j\}}
\bigl(\omega+dd^c\max\{u,-j\}\bigr)^n,
\]
where the measures for the bounded truncations are taken in the Bedford--Taylor sense. Interpret the source notation \(\mathcal E(X,\omega)\) as the full Monge--Amp\`ere mass class
\[
\mathcal E(X,\omega)=\left\{u\text{ is }\omega\text{-plurisubharmonic}:\ 
\int_X\left\langle(\omega+dd^c u)^n\right\rangle
=\int_X\omega^n\right\}.
\]
\textbf{Question (DiNezza--Lu, as attributed in the supplied record).} For every such \((X,\omega)\) and \(D\), every integrable density \(f\in L^1(X,\omega^n)\) whose restriction to \(X\setminus D\) is smooth, and every \(\varphi\in\mathcal E(X,\omega)\) satisfying
\[
\left\langle(\omega+dd^c\varphi)^n\right\rangle=f\,\omega^n
\]
as an equality of measures on \(X\), must one have
\[
\varphi|_{X\setminus D}\in C^\infty(X\setminus D)?
\]
The equation and the full mass condition imply \(f\geq0\) almost everywhere and \(\int_X f\,\omega^n=\int_X\omega^n\). The supplied statement imposes no strict positivity condition on \(f\), even on \(X\setminus D\); this absence is retained. The full mass interpretation and the normalization of \(dd^c\) above are explicit conventions, since the supplied excerpt does not define them.
\subsection{Short English statement}
On a compact K\"ahler manifold, suppose a full mass \(\omega\)-plurisubharmonic potential solves the complex Monge--Amp\`ere equation with an integrable density that is smooth away from a divisor. Must the potential also be smooth away from that divisor? The supplied wording includes densities that may vanish there.
\subsection{Sources}
\begin{itemize}
\item Supplied local UnsolvedMath registry record, permanent ID \(800890\), exact source problem code \texttt{catalogue entry}, in the AMR Open Problem Lists collection. The record labels the problem \emph{open}; this is a supplied status, not an independently verified assessment.
\item Dinew, Guedj, and Zeriahi, \emph{Open Problems in Pluripotential Theory} (2015), Question 21. Supplied source URL: \url{https://arxiv.org/abs/1511.00705}. The record attributes the question/context to DiNezza--Lu.
\item Supplied provenance: extraction method \texttt{arxiv\_latex}; source access date \texttt{2026-10-16}; upstream revision \texttt{d371ee917ac9647fbb06cd6fc529deec6f3dda82}. The access date and provenance are supplied registry information; no independent source visit or literature verification is asserted here.
\end{itemize}
\end{document}
